% Generated from ../chapters/21-a-practitioner-s-framework.md
\chapter{A Practitioner's Framework}

This chapter is not a protocol. It is a way to reason about practice.

\section{Start with constraints}

Ask what prevents the organism from adapting:

\begin{itemize}
\tightlist
\item
  pain or untreated disease;
\item
  sleep deprivation;
\item
  chronic overload;
\item
  fear and guarding;
\item
  nutritional insufficiency;
\item
  repetitive movement;
\item
  social instability;
\item
  rigid identity;
\item
  lack of meaningful challenge.
\end{itemize}

Removing a limiting constraint may produce more change than adding another optimization.

\section{Train five capacities}

\subsection{1. Structural capacity}

Strength, mobility, balance, coordination, and cardiovascular reserve provide material freedom. Train progressively and specifically.

\subsection{2. Regulatory capacity}

Practice transitions: activation to calm, focus to rest, effort to release. Meditation, breathing, slow movement, and biofeedback may help, but outcomes matter more than labels.

\subsection{3. Perceptual resolution}

Develop the ability to notice small changes without immediately interpreting them. Separate sensation, observation, and explanation.

\subsection{4. Cognitive adaptability}

Undertake difficult, ambiguous work that forces models and beliefs to evolve. Protect recovery so challenge remains adaptive.

\subsection{5. Social and existential coherence}

Relationships, responsibility, contribution, and meaning shape daily physiology through behaviour and regulation. Longevity detached from purpose can become another source of threat.

\section{Use two loops}

The \textbf{fast loop} tracks daily state: sleep, pain, mood, readiness, training, food, and practice.

The \textbf{slow loop} tracks monthly or yearly structure: strength, gait, body composition, imaging when indicated, laboratory values, cognition, work capacity, and quality of life.

Do not infer structural regeneration from a dramatic session. Look for durable changes in the slow loop.

\section{Red flags}

Stop and seek professional assessment for neurological deficits, new weakness, loss of coordination, severe or unusual headache, fainting, visual changes, chest pain, unexplained weight loss, persistent joint swelling, or worsening pain.

Do not deliberately force neck or jaw cracking. Audible events are not a performance metric.

\section{The core practice}

The deepest practical instruction is simple:

\begin{quote}
Apply enough challenge to require adaptation, enough awareness to detect compensation, enough relaxation to release obsolete control, and enough recovery to consolidate a better solution.
\end{quote}

\begin{center}\rule{0.5\linewidth}{0.5pt}\end{center}
